\section{Operator implementation and admissible outputs}\label{sec:architecture}
The refinement kernels can be implemented by Fourier, integral, or attention-based neural operators. A generic function-space layer is
\begin{equation}\label{eq:operator-layer}
 v_{\ell+1}(x)=\sigma\left(W_\ell v_\ell(x)
       +\int_D k_{\theta,\ell}(x,y)v_\ell(y)\dd y+b_\ell(x)\right).
\end{equation}
For detail refinement, one may use a shared parameter vector and set
\begin{equation}\label{eq:band-velocity}
 \widehat w_{\theta,j}(r,z;x,c)
 =\Delta_j\mathcal V_\theta(r,x+z,P_{j-1}x,j,c).
\end{equation}
The band index or its scale embedding specifies the requested refinement. Projection by $\Delta_j$ ensures the output velocity belongs to the new band. At finite resolution this can be imposed by coefficient masking in an orthonormal basis. Earlier coefficients are copied rather than regenerated.

Parameter sharing across levels is compatible with the construction. The existence theorem, however, supplies a measurable random-field law, not uniform continuity of every sample map nor universal approximation by one fixed finite network. A particular architecture must separately provide bounds on approximation, conditional context sensitivity, and numerical consistency. The summability condition on $L_j$ is sufficient and can be restrictive; energy-normalized bands can have large sensitivities even when unnormalized details decay.

\subsection{A spectral neural architecture with explicit theorem inputs}\label{sec:spectral-network}
A concrete architecture makes the abstract kernel constants computable. Index scalar orthonormal modes by $j\ge1$ and generate the first coefficient exactly. For $j\ge2$, let $R_j:H_{j-1}\to\R^q$ be linear with $\|R_j\|\le r_j$, and set
\begin{align}\label{eq:neural-detail}
 f_\theta(v)&=\sum_{k=1}^w p_k\tanh(a_k^Tv),
 \qquad p_k\ge0,\quad\sum_kp_k=1,\\
 \widehat U_j&=s_j\{f_\theta(R_j\widehat X_{j-1})+\sigma_\theta\xi_j\},
 \qquad s_j\ge0,\quad \xi_j\overset{\mathrm{iid}}\sim N(0,1).\notag
\end{align}
The weights are shared across modes; $s_j$ supplies the spectral decay. This is a coefficient-space neural generative operator: its finite outputs are restrictions of one recursively defined field, rather than separate networks indexed by array size. Its finite neural class is deliberately specified; no universality assertion is needed for the following result.

\begin{theorem}[Certified shared spectral generator]\label{thm:spectral-network}
Suppose $\sum_{j\ge2}s_j^2<\infty$, $0\le\sigma_\theta<\infty$, and the first coefficient has finite second moment. Every parameter vector in \eqref{eq:neural-detail} defines a pathwise consistent law in $\Ptwo(H)$, with
\begin{equation}\label{eq:neural-energy}
 \E\|\widehat U\|^2\le\E|U_1|^2+
 (1+\sigma_\theta^2)\sum_{j\ge2}s_j^2.
\end{equation}
Define $K_\theta=\sum_kp_k\|a_k\|$. The implemented kernels satisfy the global context bound
\begin{equation}\label{eq:neural-lip}
 L_j\le s_jr_jK_\theta.
\end{equation}
If the target detail law is
$K_j(x)=N(s_jm_j(x),s_j^2a_j^2)$ with $a_j\ge0$, then its exact fitting input is
\begin{equation}\label{eq:neural-fit}
 \eps_j^2=s_j^2\left\{\E\left[m_j(X_{j-1})-
 f_\theta(R_jX_{j-1})\right]^2+(a_j-\sigma_\theta)^2\right\}.
\end{equation}
Thus \Cref{thm:certificate} applies directly with \eqref{eq:neural-lip}--\eqref{eq:neural-fit}; if $\sum_js_j^2r_j^2<\infty$ and $\sum_j\eps_j^2<\infty$, its infinite accuracy estimate also applies.
\end{theorem}
\begin{proof}
The convex combination obeys $|f_\theta|\le1$. Independence and centering of each fresh $\xi_j$ give
$\E[\widehat U_j^2\mid\widehat X_{j-1}]
=s_j^2(f_\theta^2+\sigma_\theta^2)\le s_j^2(1+\sigma_\theta^2)$.
Orthogonality and monotone convergence prove \eqref{eq:neural-energy} and the almost-sure and mean-square convergence in \Cref{thm:representation}.
Since $\tanh$ is $1$-Lipschitz, $f_\theta$ is $K_\theta$-Lipschitz. Coupling two kernels with the same Gaussian innovation proves \eqref{eq:neural-lip}. The one-dimensional Gaussian transport formula gives \eqref{eq:neural-fit}; it remains valid when a variance or $s_j$ vanishes. Finally insert these quantities in \Cref{thm:certificate}.
\end{proof}

In \Cref{sec:learned-experiment}, $s_j=1/j$ and $R_jx=(x_1,x_{j-1})$. Here $r_2=\sqrt2$ because the first coordinate is duplicated, while $r_j=1$ for $j\ge3$. Constraining each row to $\|a_k\|\le1.3$ gives $K_\theta\le1.3$ for every training iterate and every generated context, including contexts outside the training sample. The learned law exists regardless of whether optimization succeeds. Accuracy is assessed separately.

\subsection{A finite-sample certificate for an independent conditional audit}\label{sec:validation}
An empirical fitting estimate must include sampling uncertainty. The next corollary uses the empirical Bernstein inequality of \citet[Theorem~4]{maurer2009}; the concentration inequality is established prior work. Its role here is to provide simultaneous inputs to the nonlinear hierarchy certificate.

Fix $K$ candidate networks using training data only. Draw $n\ge2$ independent validation fields from the target law, independently of training. For every candidate and mode $2\le j\le m$, suppose the audit can evaluate
\[
 Z_{i,j}=\left[m_j(X^{(i)}_{j-1})-
 f_\theta(R_jX^{(i)}_{j-1})\right]^2\in[0,B_j].
\]
This requires a known conditional mean or a justified conditional-law oracle; noisy coefficients alone do not reveal this loss. Let $\bar Z_j$ and $V_j$ be its sample mean and unbiased sample variance. For $0<\delta<1$, put $q=\log(2K(m-1)/\delta)$ and
\begin{equation}\label{eq:validation-input}
 u_j=\min\left\{B_j,\bar Z_j+
 \sqrt{\frac{2V_jq}{n}}+\frac{7B_jq}{3(n-1)}\right\},\qquad
 e_j^2=s_j^2\{u_j+(a_j-\sigma_\theta)^2\}.
\end{equation}

\begin{corollary}[Simultaneous validated law bounds]\label{cor:validated-law}
With probability at least $1-\delta$, all $K$ candidates simultaneously satisfy, for every $k\le m$,
\begin{equation}\label{eq:validated-law}
 W_2^2(\mu,\widehat\mu_k)\le\tau_k^2+\bar d_k^2,
 \qquad
 \bar d_1=0,\quad
 \bar d_j^2=\bar d_{j-1}^2+(e_j+L_j\bar d_{j-1})^2.
\end{equation}
The first coefficient is exact as above. An audited first-band error can instead initialize the same recurrence. Validation data may be shared across modes and candidates; independence is required across validation fields and from the fitted candidates.
\end{corollary}
\begin{proof}
Condition on all training data and candidate parameters. Apply the bounded empirical Bernstein inequality to $Z_{i,j}/B_j$, with failure probability $\delta/[K(m-1)]$, omitting deterministic zero-range variables. A union bound gives $\E Z_{i,j}\le u_j$ for every candidate and band. Equations \eqref{eq:neural-fit} and \eqref{eq:validation-input} then imply $\eps_j\le e_j$. Monotonicity of the nonnegative recurrence and \Cref{thm:certificate} prove \eqref{eq:validated-law}. Averaging over the independent training data preserves the probability statement.
\end{proof}

This is an a posteriori population certificate, not an optimization or minimax learning-rate theorem. It can validate modes beyond the training cutoff when their target conditional laws are available for independent auditing. For unaudited modes one must use a justified envelope, such as $\E Z_j\le B_j$, or withhold an accuracy claim there. Equation \eqref{eq:neural-energy} still establishes that the infinite learned field is well defined.


\subsection{Attention, quadrature, and norms}
For bounded fields, a normalized integral attention operator has the form
\begin{equation}\label{eq:attention}
 \mathcal A[v](x)=
 \frac{\int_D\exp(\ip{q(x,v(x))}{k(y,v(y))}/\sqrt d)\,Vv(y)\dd y}
      {\int_D\exp(\ip{q(x,v(x))}{k(y,v(y))}/\sqrt d)\dd y}.
\end{equation}
Its quadrature realization uses weights $w_i$ in both numerator and denominator. On an irregular grid, ordinary unweighted softmax integrates against the empirical sampling measure; it need not approximate the intended physical measure. Continuum attention and its discretizations are studied by \citet{calvello2024}. Likewise, a grid loss approximating an $L^2$ norm must use consistent quadrature:
\begin{equation}
 \norm{v-u}_{L^2(D)}^2\approx\sum_iw_i\norm{v(x_i)-u(x_i)}^2.
\end{equation}
A change in node count should not silently change the metric defining the training certificate.

\subsection{Convex constraints and frozen coarse coordinates}
\begin{proposition}[Projection onto admissible sets]\label{prop:constraints}
Let $\mathcal C\subset H$ be nonempty, closed, and convex, and suppose $\mu(\mathcal C)=1$. Its Hilbert-space metric projection satisfies
\begin{equation}
 W_2((\Pi_{\mathcal C})\push\eta,\mu)\le W_2(\eta,\mu),
 \qquad \eta\in\Ptwo(H).
\end{equation}
At a fixed coarse context $x$, the same statement applies to a nonempty closed convex detail set $\mathcal C_j(x)\subset E_j$ containing the support of $K_j(x,\cdot)$.
\end{proposition}
\begin{proof}
Metric projection onto a closed convex set is nonexpansive and fixes every point of the set. Project the generated component of any coupling with the target and compare squared distances; then take infima. Apply the same argument in $E_j$ for the conditional statement.
\end{proof}
Global constraint projection can modify previously generated coarse coordinates. To retain pathwise consistency, project only the new detail in a feasible fiber $\mathcal C_j(x)$. Such a fiber must be nonempty at generated contexts. Its dependence on $x$ can change the context Lipschitz constant, which must be checked for the resulting kernel. Positivity or affine conservation may admit convex formulations. Integer queue sizes, matching priorities, and event-valid paths require additional discrete mechanisms and are not covered merely by this convex projection result.

\subsection{A concrete training and sampling procedure}
\begin{enumerate}
\item Choose the state norm, orthogonal hierarchy, external conditioning information, and a reference law $\nu_j$ for each detail band.
\item From each training field $U$, form $x=P_{j-1}U$ and $b=\Delta_jU$. Draw independent $z\sim\nu_j$ and $r\sim\mathrm{Unif}(0,1)$.
\item Regress $\widehat w_{\theta,j}(r,(1-r)z+rb;x,c)$ against $b-z$ using the appropriate band norm. Regularize or constrain state and context sensitivities when useful.
\item At sampling time, integrate the detail flow conditional on the already generated coarse state, and append the resulting detail without modifying earlier bands.
\item Evaluate projected distribution errors, conditional errors on held-out paired resolutions, and physical or discrete admissibility in the metrics actually claimed. Apply the certificate only with justified population-error and stability inputs.
\end{enumerate}
If training data contain only $P_mU$, the detail pairs for levels above $m$ are unavailable. Parameter sharing can encode an extrapolation assumption, but it does not remove the identification barrier of \Cref{thm:minimax}.
